\chapter{Introduction to Cryptography}
\label{ch:introduction_to_cryptography}

Cryptography, historically referred to as cryptology, is fundamentally the study of mathematical techniques and protocols designed to enable secure communication and data storage in the presence of malicious adversaries. In modern computer security, cryptography provides the foundational primitives upon which secure systems are built.

\section{The Core Goals of Cryptography}

The application of cryptographic algorithms is intended to provide one or more of the following security features:
\begin{itemize}
    \item \textbf{Confidentiality}: Ensuring that data is accessible only to chosen, authorized entities. It prevents unauthorized parties from understanding the information even if they can intercept it.
    \item \textbf{Integrity and Freshness}: Detecting and preventing unauthorized tampering with data. Freshness is a temporal extension of integrity, preventing attackers from recording a legitimate, secure message and replaying it later (a replay attack).
    \item \textbf{Authenticity}: Certifying the origin of the data. It provides assurance that the data was indeed generated by the claimed sender.
    \item \textbf{Non-repudiation}: Establishing a mechanism where the creator of a piece of data cannot subsequently repudiate (deny) having created or sent it. This is vital for digital contracts and financial transactions.
    \item \textbf{Advanced Features}: Modern cryptography also explores complex functionalities such as proofs of knowledge or computation (e.g., zero-knowledge proofs), where a party can prove a statement is true without revealing the underlying data.
\end{itemize}

\section{A Brief Historical Perspective}

Cryptography is as old as written communication itself, originally born out of commercial needs (such as protecting the recipe for lacquer on clay tablets) and military applications (such as the Spartan scytale). For millennia, cryptographic algorithms were designed by humans to be computed by hand, using pen and paper.

\subsection{The Ancient View: A Battle of Wits}
Historically, cryptography was perceived as a continuous battle of wits between cryptographers, who ideated secret methods to obfuscate text, and cryptanalysts, who attempted to figure out the method and break the "cipher." A significant paradigm shift occurred in 1553 when Giovan Battista Bellaso separated the encryption method from the key. This meant the security of the message began to rely on a secret parameter rather than an obscure algorithm.

This concept was formalized centuries later by Auguste Kerckhoffs in 1883, who established six principles for a good cipher. The most famous, Kerckhoffs's Principle, states that a cryptographic system should be secure even if everything about the system, except the key, is public knowledge. Other principles emphasized practicality, such as the system being portable, applicable to telegraphic communication, and not imposing an excessive mental load on the operator.

\subsection{The Modern Era: From Machines to Mathematics}
The advent of mechanical computation drastically changed cryptography. The first rotor machine was invented by Edward Hebern in 1917, and the design was popularized during World War II by the German Enigma machine. The decisive cryptanalytic efforts at Bletchley Park, led by figures like Alan Turing, underscored the importance of algorithmic analysis and computation in breaking ciphers.

Following the war, cryptography transitioned from an ad-hoc battle of wits to a rigorous mathematical discipline. In 1949, Claude Shannon published "Communication Theory of Secrecy Systems," proving that mathematically unbreakable ciphers do exist. In 1955, John Nash argued that computationally secure ciphers are practically sufficient. Nash hypothesized that if the parts of a key interact complexly enough during encryption, the computational gap between the legitimate user (who takes polynomial time to encrypt/decrypt) and the attacker (who faces exponential time to break it) would become insurmountable.

\section{Fundamental Definitions and Information Theory}

Before delving into modern cryptographic constructions, we must formally define the components of a cryptosystem and understand how to measure information and uncertainty.

\subsection{Cryptographic Primitives and Spaces}
A cryptosystem operates over three primary sets:
\begin{itemize}
    \item \textbf{Plaintext space $\mathbf{P}$}: The set of all possible messages. In modern times, this is a set of bitstrings $\{0,1\}^l$.
    \item \textbf{Ciphertext space $\mathbf{C}$}: The set of all possible encrypted messages, $\text{ctx} \in \mathbf{C}$. Typically $\{0,1\}^{l'}$, where $l'$ may be greater than or equal to $l$.
    \item \textbf{Key space $\mathbf{K}$}: The set of all possible keys. Keys are usually binary strings $\{0,1\}^\lambda$, sometimes parsed into specific formats.
\end{itemize}

The system defines two deterministic functions:
\begin{itemize}
    \item \textbf{Encryption Function $\mathbb{E}$}: $\mathbf{P} \times \mathbf{K} \rightarrow \mathbf{C}$
    \item \textbf{Decryption Function $\mathbb{D}$}: $\mathbf{C} \times \mathbf{K} \rightarrow \mathbf{P}$
\end{itemize}
A fundamental requirement is \textbf{Correctness}: For all plaintexts $\text{ptx} \in \mathbf{P}$, there exist keys $k, k' \in \mathbf{K}$ such that $\mathbb{D}(\mathbb{E}(\text{ptx}, k), k') = \text{ptx}$. In symmetric cryptography, $k = k'$.

\subsection{Information Theory and Entropy}
To formally quantify security, we turn to Claude Shannon's Information Theory, which provides a mathematical framework for communication. In cryptography, we use information theory to quantitatively frame the concepts of "luck" and "guessing."

When a receiver obtains information through a channel, they are initially uncertain about the symbol. We model the sender's output as a random variable $\mathcal{X}$. Acquiring information is equivalent to losing uncertainty about the outcome of $\mathcal{X}$. A crucial point is that randomness characterizes a generative process, not an outcome; stating "00101 is a random string" is technically meaningless without context.

To measure uncertainty, we define \textbf{Entropy} $H(\mathcal{X})$. For a discrete random variable $\mathcal{X}$ with $n$ outcomes $\{x_0, \dots, x_{n-1}\}$ and probabilities $\text{Pr}(\mathcal{X} = x_i) = p_i$:
\[
H(\mathcal{X}) = \sum_{i=0}^{n-1} -p_i \log_b(p_i)
\]
When $b=2$, entropy is measured in bits. Entropy has desirable properties: it is non-negative, and combining independent uncertainties maps nicely to adding their entropies. 

Shannon's \textbf{Noiseless Coding Theorem} states that it is possible to encode the outcomes of $n$ independent and identically distributed random variables, each with entropy $H(\mathcal{X})$, into no less than $n H(\mathcal{X})$ bits per outcome without losing information. A profound consequence for cryptography is that arbitrarily compressing bitstrings without loss is impossible. Thus, cryptographic hash functions (which map arbitrary length inputs to fixed length outputs) \textit{must} discard some information. Furthermore, guessing an outcome of $\mathcal{X}$ is theoretically as hard as guessing a randomly generated bitstring of length $H(\mathcal{X})$.

However, standard Shannon entropy can sometimes mislead when evaluating cryptographic strength because highly biased distributions might yield a high average entropy while the most likely outcome is actually easy to guess. To address this, we use \textbf{Min-Entropy}:
\[
H_\infty(\mathcal{X}) = -\log_2(\max_i p_i)
\]
Min-entropy quantifies the difficulty of guessing the \textit{most common} output. If an attacker only needs to succeed once (e.g., guessing a weak key), min-entropy accurately reflects the effort required: guessing the most common outcome of $\mathcal{X}$ is exactly as hard as guessing a purely random string of length $H_\infty(\mathcal{X})$.

\begin{tcolorbox}[colback=blue!5!white,colframe=blue!75!black,title=Exam Insight]
Expect exercises calculating password or PIN entropy and success probabilities for brute-forcing attacks. You may need to estimate the success probability given an energy budget or determine the impact of hash truncation. Be prepared to identify whether standard or min-entropy is applicable based on whether the password selection is uniform or highly biased.
\end{tcolorbox}

\section{Confidentiality and Perfectly Secure Ciphers}

The primary goal of encryption is to provide confidentiality: preventing unauthorized entities from understanding the data. We must consider different \textbf{attacker models}:
\begin{itemize}
    \item \textbf{Ciphertext-only attack (CoA)}: The attacker only eavesdrops on the ciphertext.
    \item \textbf{Known-plaintext attack}: The attacker knows some plaintext-ciphertext pairs.
    \item \textbf{Chosen-plaintext attack (CPA)}: The attacker can choose arbitrary plaintexts and obtain their corresponding ciphertexts.
    \item \textbf{Chosen-ciphertext attack (CCA)}: The attacker can tamper with ciphertexts and observe how a decryption oracle reacts (up to seeing the actual decrypted value).
\end{itemize}

\subsection{Perfect Secrecy}
Shannon defined a \textbf{perfect cipher} as one where the ciphertext reveals absolutely no information about the plaintext. Formally, for all $\text{ptx} \in \mathbf{P}$ and $\text{ctx} \in \mathbf{C}$:
\[
\text{Pr}(\text{ptx sent} = \text{ptx}) = \text{Pr}(\text{ptx sent} = \text{ptx} \mid \text{ctx sent} = \text{ctx})
\]
This means observing the ciphertext $\text{ctx}$ does not alter the \textit{a priori} probability of any plaintext.

Shannon's Theorem (1949) proves that a symmetric cipher with $|\mathbf{P}| = |\mathbf{K}| = |\mathbf{C}|$ is perfectly secure if and only if:
\begin{enumerate}
    \item Every key is chosen with a uniform probability $\frac{1}{|\mathbf{K}|}$.
    \item For every plaintext and ciphertext pair, there is exactly one unique key that maps the plaintext to the ciphertext.
\end{enumerate}

A simple working example is the One-Time Pad (the Vernam Cipher). If $\mathbf{P}, \mathbf{K}, \mathbf{C}$ are binary strings of equal length, drawing a uniformly random, fresh key $k$ and computing $\text{ctx} = \text{ptx} \oplus k$ yields a perfect cipher. 

\subsection{The Impossibility of Perfect Usability}
While perfectly secure ciphers exist, they are perfectly unusable in most modern contexts. The fundamental flaw is key management. Because the key must be entirely random, never reused, and at least as long as the message itself, securely distributing and storing these keys is a logistical nightmare. Perfect ciphers implemented in history were often broken not by cryptanalysis, but through operational failures: key theft, key reuse, or flawed random key generation.

\section{Computationally Secure Cryptography}

Since perfect secrecy is impractical, modern cryptography embraces a pragmatic compromise: \textbf{Computational Security}. 

The underlying assumption is that an attacker possesses limited computational resources (bounded by polynomial time). We design ciphers such that breaking them—retrieving the plaintext without the key—is equivalent to solving a known "hard" computational problem. Examples of such assumed hard problems include factoring large integers, computing discrete logarithms, solving non-linear Boolean simultaneous equations, or finding the shortest vector in a lattice.

To formally prove computational security, cryptographers follow a rigorous paradigm:
\begin{enumerate}
    \item Define the ideal, permissible behavior of an attacker.
    \item Assume a specific computational problem is hard.
    \item Mathematically prove that if a non-ideal attacker exists who can break the cipher's security property faster than random guessing, that attacker can be used as a subroutine to efficiently solve the underlying hard computational problem.
\end{enumerate}
This shifts the burden of proof. A cipher is deemed secure as long as the underlying mathematical problem remains unsolved by efficient algorithms. An open holy grail in cryptography is to prove an exponential lower bound for these hard problems, which would shift cryptography from relying on assumptions to relying on mathematical theorems (and consequently prove $P \neq NP$).

\subsection{Security Margins and Feasibility}
When we say a problem is computationally unfeasible, we must quantify it. We evaluate security in terms of Boolean operations required to break the cipher, and we can compare this to thermodynamic energy limits (e.g., Landauer's principle). 
\begin{itemize}
    \item $2^{65}$ operations: Requires energy equivalent to boiling an Olympic swimming pool.
    \item $2^{80}$ operations: Equivalent to boiling the annual rainfall on the Netherlands. Considered legacy-level security today.
    \item $2^{128}$ operations: Considered secure for the next 5 to 10 years.
    \item $2^{256}$ operations: Provides long-term, essentially eternal security against classical computers.
\end{itemize}

A critical cautionary note: One must not blindly compare key bit-sizes across different algorithms. While brute-forcing a 128-bit symmetric key takes $\mathcal{O}(2^{128})$ operations, breaking a 2048-bit RSA key using the General Number Field Sieve takes sub-exponential time. Comparing the raw bit-size of an RSA key with an AES key is fundamentally incorrect; one must compare the actual complexity of the best known attacks.

\section{Symmetric Cryptographic Primitives}

To build practical symmetric ciphers, we rely on generating pseudo-randomness from a smaller, manageable secret key.

\subsection{CSPRNGs, PRFs, and PRPs}
A \textbf{Cryptographically Secure Pseudorandom Number Generator (CSPRNG)} is a deterministic function that takes a short, random seed (the key) of length $\lambda$ and expands it into a much longer string of length $\lambda + l$. Crucially, the output of a CSPRNG must be computationally indistinguishable from a truly uniform random string to any attacker restricted to polynomial time.

While CSPRNGs can be built from scratch, in practice they are often constructed using \textbf{Pseudorandom Permutations (PRPs)}, which are in turn conceptually related to \textbf{Pseudorandom Functions (PRFs)}.
\begin{itemize}
    \item \textbf{Random Function (RF)}: A function chosen uniformly at random from the set of all possible functions mapping inputs to outputs.
    \item \textbf{Pseudorandom Function (PRF)}: A deterministic function parameterized by a secret seed. To a polynomial-time observer who does not know the seed, the PRF's outputs cannot be distinguished from the outputs of a true RF.
    \item \textbf{Pseudorandom Permutation (PRP)}: A bijective PRF mapping an input space to itself. A concrete instantiation of a PRP is historically called a \textbf{Block Cipher}.
\end{itemize}

No formal proof exists that true PRPs exist (as it would imply $P \neq NP$). In reality, we rely on public scrutiny. Algorithms like the Advanced Encryption Standard (AES) are selected through open, multi-year public contests. They iterate a small bijective Boolean function multiple times (rounds) to thoroughly mix the input and the key. 

\subsection{Modes of Operation}
A block cipher like AES only encrypts a fixed-size block (e.g., 128 bits). To encrypt larger, arbitrary-length plaintexts, we must use a \textbf{Mode of Operation}.

\textbf{Electronic CodeBook (ECB)}: The naive approach is to split the plaintext into blocks and encrypt each block independently with the same key. This is dangerously insecure. ECB is deterministic; identical plaintext blocks encrypt to identical ciphertext blocks, perfectly preserving data patterns (such as the visual structure of an image).

\textbf{Counter (CTR) Mode}: To achieve true security, encryption must be non-deterministic (randomized) so that identical plaintexts yield different ciphertexts. CTR mode achieves this by turning a block cipher into a CSPRNG stream cipher. It encrypts a sequence of distinct values—typically a public Number Used Once (\textbf{Nonce}) concatenated with an incrementing counter. The block cipher's outputs create a pseudorandom stream that is XORed with the plaintext. Because a unique counter/nonce is used for every block, the encryption behaves like a One-Time Pad, and the cipher achieves security against Chosen Plaintext Attacks (CPA).

\begin{tcolorbox}[colback=blue!5!white,colframe=blue!75!black,title=Exam Insight]
A common exam question involves identifying information leakage in poor implementations, such as using CTR mode with predictable or reused nonces. Be prepared to analyze or prove the security/integrity of custom cryptographic protocols or encryption schemes. You will often need to find vulnerabilities, explain the corresponding exploits, and propose specific fixes or better cryptographic primitives to secure the flawed design.
\end{tcolorbox}

\subsection{Symmetric Ratcheting}
To further protect data, particularly in long-lived sessions, systems employ symmetric ratcheting. A key is passed through a one-way function (like a hash or PRNG) to generate a new key for the next message, and the old key is immediately deleted. Because the function is one-way, an attacker compromising the current state cannot roll back the state to decrypt past messages, providing a property known as Forward Secrecy.

\section{Data Integrity and Message Authentication}

Confidentiality does not imply integrity. In fact, many ciphers (especially stream ciphers like CTR mode) are highly \textbf{malleable}. An active attacker who does not know the key can flip specific bits in the ciphertext, which will predictably flip the exact same bits in the decrypted plaintext. To prevent attackers from selectively altering data, we must explicitly add integrity checks.

\subsection{Message Authentication Codes (MAC)}
A MAC provides data integrity and origin authentication. It consists of two functions: a generation function that takes a message and a secret key to produce a small \textbf{tag}, and a verification function that checks if the tag is valid for that message and key. 

Because both the sender and the receiver must share the same secret key, MACs cannot provide non-repudiation. The receiver, possessing the key, could easily forge a valid tag themselves. A widely used method for constructing a MAC using a block cipher is the CBC-MAC. Another field-proven method is HMAC, which safely combines a message and a secret key using a cryptographic hash function.

\begin{tcolorbox}[colback=blue!5!white,colframe=blue!75!black,title=Exam Insight]
You may be asked to design secure cryptographic protocols that guarantee both confidentiality and integrity (e.g., for firmware updates or a digitally controlled safe). Always consider whether the scheme prevents tampering using MACs or digital signatures. When designing these, clearly state your assumptions about key distribution, and detail every step of the communication, ensuring you explicitly include nonces or timestamps to thwart replay attacks.
\end{tcolorbox}

\subsection{Cryptographic Hash Functions}
A cryptographic hash function $H$ maps arbitrary-length bitstrings to a fixed-length output called a \textbf{digest}. To be secure, it must satisfy three computationally hard problems:
\begin{enumerate}
    \item \textbf{Preimage Resistance}: Given a digest $d$, it is unfeasible to find any input $s$ such that $H(s) = d$.
    \item \textbf{Second Preimage Resistance}: Given an input $s$, it is unfeasible to find a different input $r$ such that $H(r) = H(s)$.
    \item \textbf{Collision Resistance}: It is unfeasible to find any two distinct inputs $r \neq s$ such that $H(r) = H(s)$.
\end{enumerate}
Due to the birthday paradox, finding a collision for a $d$-bit hash takes only about $\mathcal{O}(2^{d/2})$ operations, which dictates that hash outputs must be quite large (e.g., 256 bits or more). Modern secure choices include SHA-2 and SHA-3. Older algorithms like MD5 and SHA-1 have been conclusively broken via collision attacks and must be avoided.

\section{Asymmetric Cryptography}

A major limitation of symmetric cryptography is the \textbf{key distribution problem}: Alice and Bob must securely exchange a secret key before they can communicate over an insecure channel.

In 1976, Diffie and Hellman introduced a game-changing solution: \textbf{Asymmetric Cryptography}. They devised a method for two parties to agree on a shared secret over a public, eavesdropped channel without prior communication. 

\subsection{Diffie-Hellman Key Agreement}
The Diffie-Hellman (DH) protocol relies on the Computational Diffie-Hellman (CDH) assumption within a finite cyclic group $\mathbf{G}$ with generator $g$. 
\begin{itemize}
    \item Alice picks a random secret integer $a$ and sends $g^a$ to Bob.
    \item Bob picks a random secret integer $b$ and sends $g^b$ to Alice.
    \item Alice computes $(g^b)^a = g^{ab}$.
    \item Bob computes $(g^a)^b = g^{ab}$.
\end{itemize}
Both arrive at the shared secret $g^{ab}$. An eavesdropper sees $g^a$ and $g^b$, but computing $g^{ab}$ from these public values is assumed to be computationally unfeasible. The best known attack involves solving the Discrete Logarithm Problem to find $a$ or $b$.

\subsection{Public Key Encryption and Hybrid Systems}
Following DH, Rivest, Shamir, and Adleman (RSA) invented the first practical \textbf{Public Key Encryption} system in 1977. In this paradigm, every user generates a mathematical pair of keys: a \textbf{Public Key} (shared openly) and a \textbf{Private Key} (kept strictly secret). Anyone can use Bob's public key to encrypt a message, but it is computationally intractable to decrypt that ciphertext without Bob's corresponding private key. 

While elegant, asymmetric cryptographic operations (involving large modular exponentiation or elliptic curve math) are 10 to 1000 times slower than symmetric operations. Therefore, modern systems employ \textbf{Hybrid Encryption} (or Key Encapsulation). Alice generates a random, ephemeral symmetric key, encrypts the bulk of her data using a fast symmetric cipher (like AES), and then encrypts only that small symmetric key using Bob's slow asymmetric public key. This achieves the best of both worlds.

\section{Digital Signatures and Public Key Infrastructure}

Asymmetric cryptography also solves the problem of data authentication without a pre-shared secret, enabling true \textbf{non-repudiation} through \textbf{Digital Signatures}. 

\subsection{Digital Signatures}
A digital signature scheme reverses the roles of the keys. Alice uses her \textit{Private Key} to generate a signature for a message. Anyone can use Alice's \textit{Public Key} to verify the signature. Because only Alice possesses the private key, no one else can forge the signature, and Alice cannot logically deny having signed it.

For performance reasons, Alice does not sign the entire document; instead, she computes a cryptographic hash of the document and signs the digest. Standardized algorithms include RSA (which uniquely uses the same core mathematical function for both encryption and signing) and DSA (Digital Signature Algorithm). 

Signatures are used extensively for authenticating documents and authenticating users during login systems (via challenge-response mechanisms, where a server asks a client to sign a random nonce). However, implementation flaws can lead to disaster, such as historical vulnerabilities where document viewers executed malicious macros embedded inside a securely signed PDF or Word document, leading to a mismatch between what the user thought they were signing and the actual data structure.

\subsection{The Public Key Binding Problem and Certificates}
A critical vulnerability in asymmetric cryptography is the Man-in-the-Middle (MITM) attack. If Alice wants to encrypt a message for Bob, she requests his public key. An active attacker, Mallory, could intercept the request and provide her own public key, claiming to be Bob. Mallory can then decrypt Alice's messages, read them, re-encrypt them with Bob's actual public key, and forward them along.

To thwart this, we must bind a public key to an identity using a \textbf{Digital Certificate} (most commonly standardized as ITU X.509). A certificate contains the user's identity (e.g., a domain name), their public key, the validity period, and a digital signature created by a trusted third party known as a \textbf{Certification Authority (CA)}.

When a browser connects to a secure website, it receives the site's certificate. The browser verifies the CA's signature on the certificate using the CA's own public key, which is pre-installed in the browser's "Trusted Root Store." This establishes a hierarchical chain of trust. The security of the entire internet relies on the integrity of the certificates residing in this trusted storage; if a rogue CA certificate is added, the trust model collapses entirely.

\begin{tcolorbox}[colback=blue!5!white,colframe=blue!75!black,title=Exam Insight]
Exam questions frequently involve analyzing PKI and certificate validation practices. You might be asked to compare graph-based vs. tree-based trust models, analyze CA setups for air-gapped devices, or discuss the security implications and operational trade-offs of certificate lifetime reduction. Always evaluate how these changes affect certificate revocation lists (CRLs) and online certificate status protocols (OCSP).
\end{tcolorbox}

\section{Future Directions}

Modern cryptography is continually evolving to address emerging threats and desires:
\begin{itemize}
    \item \textbf{Quantum Computing}: Large-scale quantum computers running Shor's algorithm will efficiently solve the factoring and discrete logarithm problems, completely breaking RSA, DH, and standard Elliptic Curve cryptography. NIST is currently standardizing Post-Quantum Cryptography (PQC) algorithms based on different mathematical foundations, such as lattice problems, which are believed to be resistant to quantum attacks.
    \item \textbf{Fully Homomorphic Encryption (FHE)}: There is immense desire to perform complex computations on encrypted data in the cloud without ever decrypting it. While FHE is mathematically possible today, it remains moderately to horribly inefficient and is an active area of optimization.
    \item \textbf{Side-Channel Attacks}: Cryptographic algorithms are mathematically secure, but physical implementations leak information through power consumption, electromagnetic radiation, or execution time. Modern cryptography requires integrating these side-channel capabilities into the formal attacker model to engineer secure hardware and software implementations.
\end{itemize}

\begin{tcolorbox}[colback=gray!5!white,colframe=gray!75!black,title=Chapter Summary]
This chapter covers the foundational concepts, history, and technical primitives of cryptography, transitioning from classical ciphers to modern computationally secure protocols. Key concepts include the core goals of cryptography (Confidentiality, Integrity, Authenticity, Non-repudiation), perfect secrecy versus computational security, symmetric cryptography (AES, CTR mode, MACs, Hashes), asymmetric cryptography (Diffie-Hellman, RSA), and Public Key Infrastructure (PKI).
\end{tcolorbox}

\begin{table}[htbp]
    \centering
    \begin{tabular}{|l|p{10cm}|}
        \hline
        \textbf{Concept} & \textbf{Description} \\
        \hline
        Symmetric Cryptography & Uses a single shared key for both encryption and decryption. Fast but requires secure key exchange. Includes AES and stream ciphers (CTR). \\
        \hline
        Asymmetric Cryptography & Uses a public/private key pair. Slower but solves the key distribution problem. Includes RSA and Diffie-Hellman. \\
        \hline
        MACs \& Hashes & Provide data integrity and origin authentication. Ensure messages are not tampered with. \\
        \hline
        PKI \& Certificates & Binds a public key to an identity to prevent Man-in-the-Middle attacks, establishing a chain of trust via CAs. \\
        \hline
    \end{tabular}
    \caption{Key Cryptographic Primitives}
    \label{tab:crypto_primitives}
\end{table}
